\begin{table}[H]
\centering
\caption{Detailed comparison between Mesa and Repast4Py implementations of the RCC tumor microenvironment model. The port from Mesa to Repast4Py enabled significant scalability improvements and 3D spatial modeling capabilities.}
\label{tab:detailed-mesa-vs-repast4py}
\begin{tabular}{@{}p{3.5cm}p{5cm}p{5cm}@{}}
\toprule
\textbf{Aspect} & \textbf{Mesa (Original)} & \textbf{Repast4Py (Ported)} \\
\midrule
\textbf{Framework Base} & Agent-based modeling framework for Python with single-process execution & Distributed agent-based modeling framework with MPI support \\
\midrule
\textbf{Parallelism Model} & Single-process only; all agents in one Python interpreter & Native MPI parallelism via \texttt{mpi4py}; distributed across ranks \\
\midrule
\textbf{Grid Architecture} & \texttt{MultiGrid} (2D); custom 3D extension required & \texttt{SharedGrid} with native 3D support and arbitrary dimensionality \\
\midrule
\textbf{Spatial Dimensions} & Limited to 2D (x, y); 3D requires custom implementation & Full 3D (x, y, z) with native neighbor queries and boundary handling \\
\midrule
\textbf{Agent Scalability} & Limited to $\sim$10,000 agents before performance degradation & Designed for $>$100,000 agents with MPI distribution \\
\midrule
\textbf{Glucose System} & None implemented; would require custom field overlay & Full 3D \texttt{GlucoseField} with diffusion, consumption, and gradients \\
\midrule
\textbf{Spatial Indexing} & Built-in \texttt{get\_neighbors()} with grid-based queries & Custom \texttt{spatial\_index} dict: \texttt{(x,y,z) → set[Agent]} for fast lookup \\
\midrule
\textbf{Border Handling} & Torus/Non-torus modes via \texttt{torus} parameter & Sticky/Periodic via \texttt{BorderType} enum \\
\midrule
\textbf{Agent Occupancy} & Single or multiple agents per cell via \texttt{MultiGrid} & Single or multiple via \texttt{OccupancyType} parameter \\
\midrule
\textbf{Agent Base Class} & \texttt{mesa.Agent(unique\_id, model)} & \texttt{repast4py.Agent(aid, agent\_type, rank)} with distributed ID management \\
\midrule
\textbf{Scheduling} & Mesa's step-based scheduler with \texttt{RandomActivation} & Repast4Py's \texttt{schedule\_runner} with MPI synchronization \\
\midrule
\textbf{Neighbor Queries} & \texttt{model.grid.get\_neighbors(pos, moore=True)} & Custom functions via \texttt{grid\_utils.py}: \texttt{get\_neighbors\_3d()} \\
\midrule
\textbf{Agent Movement} & \texttt{model.grid.move\_agent(agent, pos)} & Custom \texttt{move\_agent()} with spatial index updates \\
\midrule
\textbf{Memory Usage} & All agents in single process memory & Distributed across MPI ranks; reduced per-rank memory footprint \\
\midrule
\textbf{Simulation Scale} & Small tissue patches ($<$50×50×10 cells) & Large tissue volumes (100×100×50+ cells) \\
\midrule
\textbf{Data Collection} & Mesa's \texttt{DataCollector} with pandas integration & Custom \texttt{Observer} class with step-wise measurements \\
\midrule
\textbf{Visualization} & Built-in Mesa visualization server & Custom Streamlit UI with 3D plotting via Plotly \\
\midrule
\textbf{Dependencies} & \texttt{mesa}, \texttt{numpy}, \texttt{pandas}, \texttt{networkx} & \texttt{repast4py}, \texttt{mpi4py}, \texttt{numpy}, \texttt{scipy} \\
\midrule
\textbf{Diffusion Systems} & No built-in support; manual implementation required & Native 3D diffusion with \texttt{scipy.ndimage.convolve} \\
\midrule
\textbf{Biochemical Fields} & Limited to discrete agent-based signaling & Continuous 3D fields (glucose, cytokines) with spatial gradients \\
\midrule
\textbf{Configuration} & Python-based parameter files & YAML configuration with hierarchical parameter sections \\
\midrule
\textbf{Performance} & $\mathcal{O}(n^2)$ neighbor queries; single-threaded & $\mathcal{O}(1)$ spatial index lookups; parallel execution \\
\bottomrule
\end{tabular}
\end{table}